\section*{\centering 摘要}
\addcontentsline{toc}{section}{摘要}

面向增强现实、数字孪生和三维内容创作的真实场景编辑，需要同时处理两件事：从有限图像中恢复可用的空间结构，并让新增内容具有完整几何、合理外观和跨视角一致性。传统三维重建通常依赖特征匹配、Structure-from-Motion、Bundle Adjustment 和 Multi-View Stereo 等多阶段优化，流程较长，也难以满足交互速度；二维扩散模型虽然语义生成能力强，但逐视图编辑容易引入几何冲突。围绕这一问题，本文完成 Visual Geometry Grounded Transformer（VGGT）与 TRELLIS 两个官方工程的独立复现，并分析二者在真实场景三维融合编辑中的互补关系。

在独立复现之外，本文设计了一套尚未实现的融合编辑流程：使用 VGGT 从教室多视图图像中估计相机参数、深度、pointmap 和置信度，以 TRELLIS 根据文本或参考图像生成三维资产，再通过尺度、旋转和平移完成资产注册，并用地面几何、背景保持和多视图重投影约束候选结果。方案还计划利用 VGGT 置信度、相对位姿漂移和深度/pointmap cycle consistency 对候选编辑进行检查。上述注册、融合、重渲染和闭环评分均属于方法设计，并未在本次实验中实现。

现阶段的实验成果只包括两个基础模型的独立复现：VGGT 已获得多视图深度、相机位姿与场景可视化结果；TRELLIS 已跑通图像/文本条件生成及 mesh、3D Gaussian、Radiance Field 解码流程。当前尚未实现连接两者的统一执行脚本，也没有可作为实验结果报告的融合场景。因此，本文在行文中区分“已完成的模型复现”和“提出但未实现的融合编辑思想”，后者作为后续继续实现的方向。

\vspace{\baselineskip}
\noindent\textbf{关键词：}VGGT；TRELLIS；三维重建；三维生成；场景编辑；多视图一致性

\newpage
